\section{Reciprocal matrices and local fields}\label{sec:matrix-fields}

Throughout this section $b$ is an actual simple rightmost zero.
The notation $d,S,H,Q$ is that of \eqref{pr:field-notation}.
The estimates are independent of the degree.

\subsection{A matrix envelope}

\begin{proposition}\label{mf:envelope}
Put $\eta=d\Re S$, $\tau=d\Im S$, and
\[
 \Delta=3\pi^2-8,\quad K_0=\frac{9\pi^2}{\Delta},\quad
 A_0=2+\frac{24\pi}{\Delta},\quad
 \gamma=\frac{40-3\pi^2}{2\Delta}.
\]
Then
\begin{equation}\label{mf:exact-envelope}
 Qd^2\le K_0+A_0\eta+\gamma\eta^2-\tau^2/2.
\end{equation}
In particular, in the normalization \eqref{pr:normalized-moments},
writing $\sigma=c+is$,
\begin{equation}\label{mf:coarse-envelope}
 q<\frac{21}{5}\mu^2+\frac{11}{2}\mu c+\frac{c^2}{4}-\frac{s^2}{2}.
\end{equation}
\end{proposition}

\begin{proof}
Translate and scale to $b=0,d=1$. Write $v_j=-1/\xi_j$ and let
$w_k=-1/u_k$, where the $u_k$ are all the other original zeros.
Set $M=I+\mathbf1\mathbf1^*$ and
$A=M^{1/2}\diag(w_k)M^{1/2}=K+iJ$, with $K,J$ Hermitian.
Since $\Re w_k\ge0$, we have $K\ge0$. The rank-one determinant
identity and differentiation of $p(z)=zp'(0)\prod_k(1+zw_k)$ give
\[
 \det(I+z\diag(w_k)M)=\frac{p'(z)}{p'(0)}
                    =\prod_j(1+zv_j).
\]
Thus the eigenvalues of $A$, including multiplicities, are the $v_j$.

For a Hermitian matrix $B$ define
$\mathcal T(B)_{ij}=i\operatorname{sgn}(i-j)B_{ij}$, with zero
diagonal. If $x$ is a unit vector and $b_k=|x_k|^2$, diagonal
unitary conjugation removes the coordinate phases and gives
\begin{equation}\label{mf:rankone-char}
 \det(tI-\mathcal T(xx^*))
 =\frac12\left\{\prod_k(t+ib_k)+\prod_k(t-ib_k)\right\}.
\end{equation}
Indeed $tI-\mathcal T(xx^*)$ is the midpoint of an upper and a
lower triangular matrix with diagonals $t+ib_k$ and $t-ib_k$;
their difference has rank one. The determinant is affine along
a rank-one line. Zero coordinates follow by continuity.
After omitting those coordinates, the positive roots in
\eqref{mf:rankone-char} are the solutions of
$\sum_k\arctan(b_k/t)=(2j-1)\pi/2$. The left side strictly
decreases from $p\pi/2$ to zero, where $p$ is the number of
nonzero coordinates. Since $\arctan u\le u$ and $\sum b_k=1$,
the $j$th positive eigenvalue is at most
$L_j=2/((2j-1)\pi)$. The spectrum is symmetric about zero.

Let $\eta=\tr K$ and $k=\tr K^2$. The trace pairing
$\tr(K\mathcal T(xx^*))=-x^*\mathcal T(K)x$, followed by
diagonalization of either sign of $\mathcal T(xx^*)$, shows that
for every $\lambda>0$,
\begin{equation}\label{mf:schur-pair}
 \|\mathcal T(K)\|
 \le\lambda\eta+\sqrt{k}
       \left(\sum_{j\ge1}(L_j-\lambda)_+^2\right)^{1/2}.
\end{equation}
Here the diagonal entries of the conjugated $K$ are nonnegative,
sum to $\eta$, and have sum of squares at most $k$. This proves
the inequality by Cauchy--Schwarz and includes the whole spectrum.

In an upper triangular Schur basis for $A$,
$J=\diag(\Im v_j)+\mathcal T(K)$, whence
$\|J\|\le1+\|\mathcal T(K)\|$. Write $w_j=a_j+ib_j$ and
$\beta=\sum b_j=(\Im S)/2$. Direct multiplication gives
\[
 k=3\sum a_j^2+\Bigl(\sum a_j\Bigr)^2,\quad
 \tr J^2=3\sum b_j^2+\beta^2,\quad
 \|Je\|^2=\frac{N+1}{N}\left(\sum b_j^2+\beta^2\right),
 \qquad e=N^{-1/2}\mathbf1.
\]
In particular,
\[
 \tr J^2=\frac{3N}{N+1}\|Je\|^2-2\beta^2
 \le3\|J\|^2-2\beta^2.
\]
In a Schur basis the off-diagonal squared contributions to
$\tr K^2$ and $\tr J^2$ agree, so
\[
 \tr K^2-\tr J^2
 =\sum_j\bigl((\Re v_j)^2-(\Im v_j)^2\bigr).
\]
Also $0\le\Re v_j\le1$, whence $\sum_j(\Re v_j)^2\le\eta$.
Therefore
\[
 Q=2\sum_j(\Re v_j)^2-k+\tr J^2
 \le2\eta+3(1+\|\mathcal T(K)\|)^2-k-2\beta^2.
\]
Take
\[
 \lambda=\frac{24\pi+(40-3\pi^2)\eta}
                   {6\pi(3\pi+4\eta)},\qquad
 a_\lambda^2=(2/\pi-\lambda)^2+(2/(3\pi)-\lambda)^2.
\]
For $\eta\ge0$, one has $L_3<\lambda<L_1$ and
$3a_\lambda^2<1$. These follow by multiplying the positive
denominators; the endpoint values of the fractional linear
function $\lambda(\eta)$ suffice. Thus the sum in
\eqref{mf:schur-pair} is at most $a_\lambda^2$, including when
$\lambda>L_2$. Maximizing the resulting quadratic in $\sqrt k$
gives
\[
 3(1+\|\mathcal T(K)\|)^2-k
 \le\frac{3(1+\lambda\eta)^2}{1-3a_\lambda^2}.
\]
Put
\[
 R_\eta=18\pi^2+48\pi\eta+(40-3\pi^2)\eta^2.
\]
Direct substitution gives
\[
 1+\lambda\eta=\frac{R_\eta}{6\pi(3\pi+4\eta)},\qquad
 1-3a_\lambda^2
 =\frac{\Delta R_\eta}{6\pi^2(3\pi+4\eta)^2}.
\]
Consequently
\[
 2\eta+\frac{3(1+\lambda\eta)^2}{1-3a_\lambda^2}
 =2\eta+\frac{R_\eta}{2\Delta}
 =K_0+A_0\eta+\gamma\eta^2.
\]
This proves \eqref{mf:exact-envelope}. Finally $\pi>157/50$
implies $K_0<21/5$, $A_0<11/2$, and $\gamma<1/4$.
Divide by $(Hd)^2$ to obtain \eqref{mf:coarse-envelope}.
\end{proof}

\subsection{Logarithmic bounds and saddle moments}

The next inequality controls the logarithm on a disk.

\begin{lemma}\label{mf:log-metric}
For $0<r<1$ put
\[
 h(r)=\frac{r-\log(1+r)}{r^2},\qquad
 j(r)=\frac{-r-\log(1-r)}{r^2}.
\]
For every $z=x+iy$ with $|z|\le r$,
\[
 \log|1+z|\le x-h(r)x^2+j(r)y^2.
\]
\end{lemma}
\begin{proof}
The analytic identity
$(z-\Log(1+z))/z^2=\int_0^1t/(1+tz)\,dt$ holds in $|z|<1$.
For fixed $t$, the image of $|z|\le r$ under $t/(1+tz)$ lies in
the disk with center $t/(1-t^2r^2)$ and radius
$t^2r/(1-t^2r^2)$. The difference of squared radius and squared
distance from the center is
\[
 \frac{t^4(r^2-x^2-y^2)}
 {(1-t^2r^2)((1+tx)^2+t^2y^2)}\ge0.
\]
Integration therefore puts the preceding integral in the disk
with center $c_r=-\log(1-r^2)/(2r^2)$ and radius
$e_r=(\operatorname{atanh}r-r)/r^2$. Multiplication by $z^2$
and taking real parts gives
$\Re\Log(1+z)\le x-(c_r-e_r)x^2+(c_r+e_r)y^2$.
The two coefficients are $h(r)$ and $j(r)$.
\end{proof}

\begin{lemma}\label{mf:saddle-dual}
Let $\Re v_j\ge0$, $|v_j|\le1$, $\eta=\sum\Re v_j$, and
$Q=\sum|v_j|^2>0$. Suppose $z_0=-X+iY$, $|z_0|=s<1$,
$X\ge0$, and $\sum_jv_j/(1+z_0v_j)=0$. Then
\[
 \frac{\eta}{Q}\ge
 \frac{X(1-s)^2}{(1+|Y|)^2-X^2}.
\]
\end{lemma}
\begin{proof}
Conjugate to make $Y\ge0$, and put
$B=X/(1+2Y+s^2)$, $A=(1-2BX)/(1-s)^2$.
For $v=a+ib$, $r=|v|\le1$, and $d_v=|1+z_0v|^2$, expansion gives
\[
 (Aa-Br^2)d_v-(a-Xr^2)
 =a(Ad_v-1+2BXr^2)
   +Br^2[2Y(1+b)+s^2(1-r^2)].
\]
The second summand is nonnegative. Also $d_v\ge(1-sr)^2$ and
$2BX\le2s^2/(1+s^2)\le s<1$. The function
$A(1-sr)^2+2BXr^2$ decreases on $[0,1]$ to its value one
at $r=1$, since $As(1-s)-2BX=(s-2BX)/(1-s)\ge0$.
Thus $\Re(v/(1+z_0v))\le Aa-B|v|^2$. Sum and use the saddle
equation to obtain $\eta/Q\ge B/A$, the stated bound.
\end{proof}

\begin{lemma}\label{mf:saddle-variance}
Let $\Re v_j\ge0$, $|v_j|\le1$, $Q=\sum|v_j|^2>0$, and
$\sum_jv_j/(1+z_0v_j)=0$, with $\Re z_0\le0$.
Write $u_j=v_j/(1+z_0v_j)$ and $Q_u=\sum|u_j|^2$.
Suppose $|z_0|<s<1$, $-\Re z_0<X$, $|\Im z_0|<Y$,
and $X<1-Y$. Here $s$ is an upper bound for the modulus.
Put $U=(1-s)^{-1}$.
Then
\[
 \frac1{1+2Y+s^2}\le\frac{Q_u}{Q}
 \le q_m:=\min\left(U^2,\frac1{1-2X^2-2Y}\right)
\]
when the last denominator is positive, and
\[
 \frac{\sum(\Re u_j)^2}{Q_u}\le\beta
 :=\min\left\{
 X\left(\frac{1+X}{1-s^2}+XU^2\right),
 \frac{X(1-Y+X)}{(1-Y)^2+X^2}\right\}.
\]
\end{lemma}
\begin{proof}
Write $z_0=x_0+iy_0$, $v_j=a_j+ib_j$, and
$d_j=|1+z_0v_j|^2$. Then
\[
 \Re u_j=\frac{a_j+x_0|v_j|^2}{d_j},\quad
 \Im u_j=\frac{b_j-y_0|v_j|^2}{d_j},\quad
 |u_j|^2=\frac{|v_j|^2}{d_j}.
\]
Since $a_j\ge0$, $x_0\le0$, and $|v_j|\le1$,
$d_j\le1+2Y+s^2$. The zero real sum gives
$\sum a_j/d_j=-x_0Q_u$. In the exact identity
\[
 Q=Q_u+2x_0\sum a_j|u_j|^2
       -2y_0\sum b_j|u_j|^2
       +|z_0|^2\sum|v_j|^2|u_j|^2
\]
we have $0\le\sum a_j|u_j|^2\le-x_0Q_u$ and
$|\sum b_j|u_j|^2|\le Q_u$. This proves the mass bounds.

The inverse disk inequality
$(1-|z_0|^2)|u_j|^2+2\Re(z_0u_j)\le1$ gives
$\Re u_j\le(1+X)/(1-s^2)$, while
$\Re u_j\ge-X|u_j|^2\ge-XU^2$.
Split the zero real sum into its positive and negative terms.
The total absolute negative part is at most $XQ_u$. Multiplying
each part by its respective amplitude bound proves the first
variance bound.

For the second, write $z_0=-X_0+iY_0$. Give
$r_j=1/\overline{v_j}$ weights $|u_j|^2/Q_u$, omitting a
zero $v_j$ if necessary. Since
$1/\overline{u_j}=r_j+\overline{z_0}$ and $\sum u_j=0$,
their weighted mean is $X_0+iY_0$. Each $r_j$ has nonnegative
real part and modulus at least one. Put $L=1-|Y_0|$ and
$c=X_0/L<1$. The numbers
$r'_j=(r_j-iY_0)/L=a_j'+ib_j'$ have $a_j'\ge0$,
$a_j'^2+b_j'^2\ge1$, and weighted mean real part $c$.
On this unbounded domain,
\[
 \frac{(a-c)^2}{(a-c)^2+b^2}\le\frac{a+c^2}{1+c^2}.
\]
For $a\ge1$ the right side is at least one. For $0\le a\le1$
replace $b^2$ by $1-a^2$ and use
\[
 (a+c^2)(1-2ac+c^2)-(1+c^2)(a-c)^2
 =a(1-a)(1+c)^2\ge0.
\]
The denominator cannot vanish. Taking weighted means proves
$\sum(\Re u_j)^2/Q_u\le c(1+c)/(1+c^2)$.
This function increases on $[0,1]$, with derivative
$(1+2c-c^2)/(1+c^2)^2>0$. Substituting $c\le X/(1-Y)$
proves the second bound.
\end{proof}

\subsection{An affine reciprocal-field estimate}

\begin{proposition}\label{mf:affine}
For every actual simple rightmost zero,
\begin{equation}\label{mf:affine-bound}
 Qd^2\le20+11Hd.
\end{equation}
Consequently $Qd^2\le13(1+Hd)$ whenever $Hd\ge7/2$.
\end{proposition}

\begin{proof}
Normalize to $b=0,d=1$, and suppose $Q>20+11H$.
Set $B=-\sum v_j^2$, $\eta=\Re S$.
The matrix envelope implies $\eta>0$. Conjugate so that
$\Im S\ge0$, and write
\[
 t=\frac{\Im S}{\eta},\quad
 x=\frac{\Re B}{Q},\quad y=\frac{|\Im B|}{Q}.
\]
We have $B-Q=-2\sum(\Re v_j)v_j$, hence
$|B-Q|\le2\eta$ and $|B|\le Q$. Also
$H/Q<1/11$. The envelope gives
$t^2<2\gamma<(7/10)^2$: if its quadratic coefficient were
nonpositive, its constant and linear coefficients would already
be smaller than $20+11H$. Thus
\begin{equation}\label{mf:affine-domain}
 (t,y)\in[0,7/10]\times[0,1/4],\quad
 (x-1)^2+y^2<\frac4{121(1+t^2)},\quad x^2+y^2\le1.
\end{equation}
Write $B/Q=x+i\epsilon y$, with $\epsilon\in\{-1,1\}$.
Both signs are used when $y=0$.

Here and below directed rational rounding is denoted by
$\lfloor u\rfloor_m=\lfloor mu\rfloor/m$ and
$\lceil u\rceil_m=\lceil mu\rceil/m$. Write
$\operatorname{sf}_m(u)=\lfloor\sqrt u\rfloor_m$ and
$\operatorname{sc}_m(u)=\lceil\sqrt u\rceil_m$; these are computed
by integer square roots. Unless indicated, square-root meshes
have $m=10^6$. Use $p_0=6283/2000<\pi<22/7$ and put
$k_U=9p_0^2/(3p_0^2-8)$,
$a_U=2+24p_0/(3p_0^2-8)$,
$g_U=(40-3p_0^2)/(2(3p_0^2-8))$.
These are upper bounds for $K_0,A_0,\gamma$, since each decreases
on the indicated interval for $\pi$.

Each row of the certificate consists of six rationals
$(t_l,t_h,y_l,y_h,\rho,k)$, with $\rho>0$, $k>0$, and the
phase sign $\epsilon$.
They specify a closed rectangle and a contour. The following
construction gives every bound used to check that row.
Put
\[
 d_l=\operatorname{sf}(1+t_l^2),\quad
 d_h=\operatorname{sf}(1+t_h^2),\quad
 e_0=\frac2{11d_l},\quad
 g_0=g_U-t_l^2/2.
\]
If $y_l>e_0$ the cell is impossible. Otherwise define
$r_0=1-\operatorname{sc}(e_0^2-y_l^2)$.
If $r_0^2+y_l^2\ge1$ or $g_0\le0$ the cell is also impossible
under \eqref{mf:affine-domain}. In a retained cell
cap $y_h$ by
$\widehat y_h=\min(y_h,e_0,\operatorname{sc}(1-r_0^2))$.
The retained rows have $r_0>\widehat y_h\ge y_l$ and $g_0>0$.
Let $h_0$ be the largest $1/100$ mesh point strictly below the
positive root of
\[
 P(h)=k_U-20+(a_U-11d_l)h+g_0h^2.
\]
The checks $h_0>0$, $P(h_0)<0$, and $P(h_0+1/100)\ge0$
imply $\eta>h_0$. Set
\[
 Q_0=20+11d_lh_0,\quad r_Q=\operatorname{sf}_{10^4}(Q_0),\quad
 b_m=\operatorname{sf}(r_0^2+y_l^2).
\]
Choose the first $10^{-6}$ mesh point $a\in[1,2]$ with
\begin{equation}\label{mf:saddle-rouche}
 (ab_m-1)(11-a)>a^2.
\end{equation}
Put $s=a/11$, $U=(1-s)^{-1}$, $e=a^2/[b_m(11-a)]$.
For $L=p''/p'$ the geometric series gives
\[
 L(z)=S+Bz+E(z),\qquad |E(z)|\le\frac{Q|z|^2}{1-|z|}.
\]
On $|z|=aH/Q<s$, \eqref{mf:saddle-rouche} gives
$|Bz|>|S+E(z)|$. Rouch\'e's theorem gives an actual zero $z_0$
of $p''$ with $|z_0|<aH/Q$ and
\[
 |z_0+S/B|<(H/Q)e,\qquad \Re z_0\le0.
\]
The latter assertion is Gauss--Lucas twice; $p'$ is nonzero
in the whole disk $|z|<1$.

For phase bounds set
\[
 q_U=1/d_l,\ c_U=t_h/d_h,\quad
 q_L=1/\operatorname{sc}(1+t_h^2),\
 c_L=t_l/\operatorname{sc}(1+t_l^2),
\]
\[
 R_f=\frac{r_0}{r_0^2+y_l^2},\quad
 Y_f=\frac{\widehat y_h}{r_0^2+\widehat y_h^2},\quad
 R_L=\operatorname{sf}(1-\widehat y_h^2).
\]
On the unit-disk fiber, differentiation in $x>y$ gives
$R_L\le x/(x^2+y^2)\le R_f$ and
$y_l\le y/(x^2+y^2)\le Y_f$. For $\epsilon=1$ take
\[
 D_c=q_UR_f+c_UY_f+e,\qquad
 I_c=\max(0,c_UR_f-q_Ly_l,q_UY_f-c_LR_L)+e;
\]
for $\epsilon=-1$ take
\[
 D_c=\max(0,q_UR_f-c_Ly_l)+e,\qquad
 I_c=c_UR_f+q_UY_f+e.
\]
These follow by multiplying $S/B$ explicitly, and give
$-\Re z_0<(H/Q)D_c$, $|\Im z_0|<(H/Q)I_c$.
Initially put $X=\min(s,D_c/11)$, $Y=\min(s,I_c/11)$.
Lemma~\ref{mf:saddle-dual} permits replacing $X$ by its minimum
with
\[
 \frac{\operatorname{sc}((1-s)^4+4h_*^2(1+Y)^2)-(1-s)^2}
      {2h_*},\qquad h_*=\frac1{11d_l}.
\]
Apply Lemma~\ref{mf:saddle-variance} and denote its upper
mass and variance bounds by $q_m,\beta$. Each row checks
$1-2X^2-2Y>0$, $X<1-Y$, and $\beta<1/2$.

For $u_j=v_j/(1+z_0v_j)$ let $B_0=-\sum u_j^2$.
The geometric series gives
$|B_0-B|/Q<U^2-1=:p_*$. The mass and variance bounds give
\[
 \Re(B_0/Q)\ge\frac{1-2\beta}{1+2Y+s^2},\quad
 |\Im(B_0/Q)|\le2q_m\sqrt{\beta(1-\beta)},\quad
 |B_0/Q|\le q_m.
\]
Use the rational enclosures
\[
 R=\left\lfloor\max\left(r_0-p_*,
             \frac{1-2\beta}{1+2Y+s^2}\right)\right\rfloor_{10^4},
\]
\[
 I=\left\lceil\min(\widehat y_h+p_*,
             2q_m\operatorname{sc}(\beta(1-\beta)))\right\rceil_{10^4},
 \quad M=\lceil q_m\rceil_{10^4},\quad
 G=\operatorname{sf}_{10^4}(2p_0/M).
\]
Thus $R>0$, and
\[
 \left|\int_{-\infty}^{\infty}e^{-B_0t^2/(2Q)}\,dt\right|
 =\sqrt{\frac{2\pi}{|B_0/Q|}}\ge G.
\]

\smallskip
\noindent\emph{The contour.}
Define
\[
 \lambda=\left\lfloor\frac{\rho}{k/r_Q+D_c/11}\right\rfloor_{10^4},
 \quad T_0=\lfloor\rho r_Q\rfloor_{10^4},\quad
 \delta=\operatorname{sc}_{10^5}(\rho^2(1+\lambda^{-2})).
\]
Check $\lambda>1$, $T_0>0$, $U\delta<1$. Put
\[
 T=\max\{T_0,\lambda[k+(H/\sqrt Q)D_c]\},\qquad
 D_c'=k-\sqrt Q\Re z_0.
\]
Then $0<D_c'<T/\lambda$ and $T/\sqrt Q<\rho$.
Indeed $Q>r_Q^2$ and $H/Q<1/11$, so
\begin{align*}
 \frac{T}{\sqrt Q}
 &=\max\left\{\frac{T_0}{\sqrt Q},
       \lambda\left(\frac{k}{\sqrt Q}+\frac{H}{Q}D_c\right)\right\}\\
 &<\max\left\{\frac{T_0}{r_Q},
       \lambda\left(\frac{k}{r_Q}+\frac{D_c}{11}\right)\right\}
 \le\rho.
\end{align*}
Take the vertical segment $z_0+it/\sqrt Q$, $-T\le t\le T$,
and both horizontal segments joining its ends to $\Re z=k/\sqrt Q$.
Relative to $z_0$ every point has modulus less than $\delta$.

Use the directed logarithm coefficients
\begin{equation}\label{mf:directed-log}
 \underline h(r)=\left\lfloor\sum_{\ell=0}^{15}
                      \frac{(-r)^\ell}{\ell+2}\right\rfloor_{10^6},
 \qquad
 \overline j(r)=\left\lceil\sum_{\ell=0}^{15}
          \frac{r^\ell}{\ell+2}+\frac{r^{16}}{18(1-r)}
                       \right\rceil_{10^6}.
\end{equation}
The alternating and geometric remainders prove
$\underline h\le h$ and $\overline j\ge j$.
Replace the two coefficients in Lemma~\ref{mf:log-metric}
before expanding its quadratic form. At $r=U\rho$ and
$r=U\delta$ write the respective pairs as $(h_v,j_v)$ and
$(h_c,j_c)$. Set
\[
 q_l=\lfloor(1+2Y+s^2)^{-1}\rfloor_{10^6},\quad
 A_v=\lfloor h_v-(h_v+j_v)\beta\rfloor_{10^6},
\]
\[
 A_c=\lfloor h_c-(h_c+j_c)\beta\rfloor_{10^6},\quad
 B_c=\lceil j_c-(h_c+j_c)\beta\rceil_{10^6}.
\]
Every row has $q_l,A_v>0$ and $B_c\ge0$.

\smallskip
\noindent\emph{The paired vertical integral.}
For $|t|\le T$ put
$f(t)=\sum_j\Log(1+iu_jt/\sqrt Q)$ and
$g(t)=-B_0t^2/(2Q)$. For a symmetric shell
$a_j\le|t|\le b_j$, put
$\rho_j=\min(\rho,b_j/r_Q)$, taking $\rho_j=\rho$ if $b_j=\infty$,
and $r_j=U\rho_j$. With $h_j=\underline h(r_j)$ and
$j_j=\overline j(r_j)$, set
\[
 \alpha_j=\left\lfloor\min\{R/2,
                q_l[h_j-(h_j+j_j)\beta]\}\right\rfloor_{10^5}>0.
\]
Both $\Re f(\pm t)$ and $\Re g(t)$ are at most $-\alpha_jt^2$.
Writing the even and odd parts of $f-g$ as $E(t),O(t)$,
the logarithmic series gives
\[
 |O(t)|\le\frac{q_mU|t|^3}{\sqrt Q}
       \left(\frac13+\frac{r_j^2}{5(1-r_j^2)}\right),\quad
 |E(t)|\le\frac{q_mU^2t^4}{Q}
       \left(\frac14+\frac{r_j^2}{6(1-r_j^2)}\right).
\]
These bounds retain all odd orders at least five and all even
orders at least six. For arbitrary complex $A,O$,
\[
 e^{A+O}+e^{A-O}-2e^A
 =O^2\int_{-1}^1(1-|s|)e^{A+sO}\,ds.
\]
Apply this with $A=g+E$, and apply the first-difference integral
to $e^{g+E}-e^g$. The real-part bounds hold along both line
segments, so
\[
 |e^{f(t)}+e^{f(-t)}-2e^{g(t)}|
 \le(|O(t)|^2+2|E(t)|)e^{-\alpha_jt^2}.
\]
Use the shell endpoints $0,1,2,3,4,5,6,8,10,12,\infty$.
If $\Delta M_{k,j}$ is an upper bound for
$\int_{a_j}^{b_j}t^ke^{-\alpha_jt^2}\,dt$, the
paired error is at most
\[
 \mathcal E=\left\lceil\sum_j
 \frac{q_m^2U^2c_{o,j}^2\Delta M_{6,j}
          +2q_mU^2c_{e,j}\Delta M_{4,j}}{r_Q^2}
                      \right\rceil_{10^4},
\]
where $c_{o,j}$ and $c_{e,j}$ are the parenthesized coefficients.
Extending a shell past the actual endpoint $T$ only enlarges its
positive error integral.

Here the moments have a purely rational verification. For
$M_k(a)=\int_a^\infty t^ke^{-\alpha t^2}\,dt$,
\[
 M_k(a)=\frac{a^{k-1}e^{-\alpha a^2}}{2\alpha}
          +\frac{k-1}{2\alpha}M_{k-2}(a).
\]
For $a>0$,
$ae^{-\alpha a^2}/(1+2\alpha a^2)\le M_0(a)
 \le e^{-\alpha a^2}/(2\alpha a)$.
Differentiating these proposed bounds and integrating to infinity
proves the two signs. At zero,
$M_4(0)=3\sqrt\pi/(8\alpha^{5/2})$ and
$M_6(0)=15\sqrt\pi/(16\alpha^{7/2})$.
Use directed square roots and $\pi<22/7$.
For each finite endpoint $x=\alpha a^2<82$,
\[
 E_{40}(x)\le e^x\le E_{80}(x)
               +\frac{x^{81}}{81!(1-x/82)},\qquad
 E_m(x)=\sum_{\ell=0}^m\frac{x^\ell}{\ell!}.
\]
Taking reciprocals in the correct directions gives rational
moment bounds; subtract a lower right tail from an upper left
tail. The final infinite shell is included. The certificate
checks every domain and positivity condition in this construction.

\smallskip
\noindent\emph{The two horizontal segments.}
There are two valid bounds, either of which may be used on a row.
For the translated bound, expansion of Lemma~\ref{mf:log-metric}
at $w=(d+i\epsilon_c T)/\sqrt Q$, with
$\epsilon_c\in\{-1,1\}$ and $0\le d\le T/\lambda$, gives
\[
 \Re\sum\Log(1+u_jw)
 \le-q_l(A_c-B_c/\lambda^2)T^2
           +(h_c+j_c)IT^2/\lambda.
\]
Indeed, if $\nu=\sum(\Re u_j)^2/Q_u\le\beta$, the quadratic
part before substitution is
\[
 \frac{Q_u}{Q}\{-[h_c-(h_c+j_c)\nu]T^2
                  +[j_c-(h_c+j_c)\nu]d^2\},
\]
which increases with $\nu$ since $T^2\ge d^2$.
The mixed term has modulus at most $(h_c+j_c)IdT$.
When $A_c-B_c/\lambda^2>0$, the downward-rounded damping is
\[
 \kappa_{\rm tr}=
 \left\lfloor q_l(A_c-B_c/\lambda^2)
                  -(h_c+j_c)I/\lambda\right\rfloor_{10^5}.
\]

For the original-coordinate bound put
$l=1/\lambda$, $J=1+I_c/(\lambda D_c)$, and
$r_o=\operatorname{sc}_{10^5}(\rho^2(l^2+J^2))$.
If $J<2$ and $r_o<1$, take $h_o=\underline h(r_o)$,
$j_o=\overline j(r_o)$, and
$A_o=h_o-(h_o+j_o)/(11d_l)$.
Writing $z=x_c+iy_c$, $V=\sum(\Re v_j)^2\le\eta$ and
$W=\sum\Re v_j\Im v_j$, the exact expanded bound is
\[
 \log|p'(z)/p'(0)|
 \le\eta x_c-(\Im S)y_c
 +[j_oQ-(h_o+j_o)V]x_c^2
 +[(h_o+j_o)V-h_oQ]y_c^2+2(h_o+j_o)Wx_cy_c.
\]
On both segments,
$|x_c|\sqrt Q/T\le l$,
$|y_c|\sqrt Q/T\in[2-J,J]$,
$x_c\le k/\sqrt Q$, and $H/(\sqrt Q T)\le1/(\lambda D_c)$.
As $2|W|\le\widehat y_hQ$, the normalized upper bound is
\[
 j_ol^2+\frac{q_Uk}{\lambda D_cT_0}
 +\max_{2-J\le v\le J}(-A_ov^2+Lv),\qquad
 L=(h_o+j_o)\widehat y_hl+\frac{c_U}{\lambda D_c}.
\]
When $A_o>0$, the maximum is attained at an endpoint or at
$L/(2A_o)$ if that point belongs to the interval.
For normalization at the saddle, its exact equation gives
\[
 \ell_0=\sum_j\left[\Log(1+z_0v_j)
                      -\frac{z_0v_j}{1+z_0v_j}\right].
\]
The series has coefficient $(-1)^m(m-1)/m$ in degree
$m\ge2$. Hence $|\ell_0|\le Q|z_0|^2f(s)$, where
$f(s)=1/(1-s)-\sum_{j\ge0}s^j/(j+2)$.
Use
$f_s=\lceil1/(1-s)-\sum_{j=0}^{15}s^j/(j+2)\rceil_{10^6}$.
Since $Q|z_0|^2/T^2\le a^2/(\lambda^2D_c^2)$, a second
valid damping is
\[
 \kappa_{\rm or}=\left\lfloor
 -\max_{[2-J,J]}(-A_ov^2+Lv)-j_ol^2
 -\frac{q_Uk}{\lambda D_cT_0}
 -\frac{f_sa^2}{\lambda^2D_c^2}\right\rfloor_{10^5}.
\]
Choose the larger valid positive damping $\kappa$.
Every retained row has at least one.

\smallskip
\noindent\emph{The final finite predicates.}
Put
$F_0=\lfloor(G-\mathcal E-1/1000-1/100)/2\rfloor_{10^4}$
and $J_0=\lceil1+I_c/(\lambda D_c)\rceil_{10^4}$.
In addition to the conditions above, each row verifies
\begin{gather}
 F_0>0,\qquad
 E_2(RT_0^2/2)>\frac{2000}{RT_0},\qquad
 \kappa T_0^2>\frac12,\label{mf:affine-final-one}\\
 E_{24}(\kappa T_0^2)>\frac{200T_0}{\lambda},\qquad
 \kappa(k^2+J_0^2T_0^2)>J_0^2,\label{mf:affine-final-two}\\
 E_{24}(\kappa T_0^2)>
       \frac{11}{10}\frac{k^2+J_0^2T_0^2}{kF_0}.
 \label{mf:affine-final-three}
\end{gather}
The two Gaussian tails are bounded by
\[
 2\int_T^\infty e^{-Rt^2/2}\,dt
 \le\frac{2e^{-RT^2/2}}{RT}<\frac1{1000}.
\]
The sum of the moduli of the two horizontal integrals, divided by
$|p'(z_0)|/\sqrt Q$, is bounded by
\[
 \frac{2T}{\lambda}e^{-\kappa T^2}<\frac1{100}.
\]
The first bound decreases for $T>0$; the second decreases for
$\kappa T^2>1/2$. Thus the finite comparisons at $T_0$ prove both
bounds on the whole ray $T\ge T_0$.
The endpoint primitive values therefore differ by more than
$2F_0|p'(z_0)|/\sqrt Q$. At one endpoint
$z_{\epsilon_c}=k/\sqrt Q+i(\Im z_0+\epsilon_c T/\sqrt Q)$,
\[
 |p(z_{\epsilon_c})|>F_0|p'(z_0)|/\sqrt Q,\quad
 |p'(z_{\epsilon_c})|\le e^{-\kappa T^2}|p'(z_0)|,\quad
 \sqrt Q|\Im z_{\epsilon_c}|\le J_0T.
\]
But every original zero lies in $\Re z\le0$, and zero itself
is a simple zero. Therefore
$\Re(p'/p)(z)\ge\Re(1/z)$ in $\Re z>0$, which would imply
\[
 \frac{k}{k^2+J_0^2T^2}<\frac{e^{-\kappa T^2}}{F_0}.
\]
The last two predicates make
$(k^2+J_0^2T^2)e^{-\kappa T^2}$ decreasing for $T\ge T_0$
and already smaller than $(10/11)kF_0$ at $T_0$.
This is a contradiction.

The finite input consists of the four tables
\cert{C11_A20_PHASE_1_DIRECT_QUADRATIC_METRIC_CELLS.csv},
\cert{C11_A20_PHASE_-1_DIRECT_QUADRATIC_METRIC_CELLS.csv},
and their corresponding \cert{SKIPS.csv} tables.
They contain respectively $77,198$ retained rectangles and
$14,14$ excluded rectangles. On each sign they cover exactly
$[0,7/10]\times[0,1/4]$. This is checked by sorting every
rational vertical edge, merging the horizontal intervals on
each intervening open slab, and repeating the check on each
closed edge. All excluded rectangles satisfy the explicit
exclusion above. The program
\cert{affine/replay_affine_c11_a20.py} reconstructs every
quantity from the six rational entries and verifies
\eqref{mf:saddle-rouche}--\eqref{mf:affine-final-three}
with rational arithmetic, including all $2750$ shells.
This proves the required finite assertion and the proposition.
The final consequence follows from $20+11D\le13(1+D)$
for $D\ge7/2$.
\end{proof}

\subsection{Closing the unbounded field range}

The two signed-curvature bounds from Proposition~\ref{lt:quartic}
give the following estimate.

\begin{proposition}\label{mf:D270}
If $D=Hd\ge270$, then $Qd/H<9/4$.
\end{proposition}
\begin{proof}
Put $y=Qd/H=qD$. Proposition~\ref{mf:affine} implies
\[
 y\le13(1+1/D)\le Y_0:=3523/270,\qquad
 q\le13(D+1)/D^2\le3523/72900<97/2000.
\]
We also have $\mu=1/D\le1/270$.
For $0\le r\le97/2000$ define
\[
 v(r)=r+17r^2+137r^3+313r^4,\quad m(r)=1-v(r),
\]
\[
 A(r)=\frac{1+17r+137r^2+313r^3}{m(r)},\qquad
 B(r)=\frac{18+154r+450r^2+313r^3}{m(r)},
\]
\[
 C(r)=1+\left(2+\frac{100}{199}\right)rA(r)^2
 +3rB(r)+\frac{600}{199}r^2A(r)^2
 +4rA(r)+29r+436r^2.
\]
Let $K=p''/p'$. The coarse quartic row gives
\[
 \frac{\Re K'(b)}{H^2}
 <C(q)q^2+\frac{50000}{3}q^5+\frac{650}{9}q\mu.
\]
On the other hand, rightmostness gives
$\Re K'(b)=Q-2\sum(\Re v_j)^2\ge Q-2H/d$.
Consequently
\begin{equation}\label{mf:coarse-bootstrap}
 0<P_D(y):=y^2C(y/D)+\frac{50000}{3}\frac{y^5}{D^3}
                       +\frac{650}{9}y-D(y-2).
\end{equation}

On the whole interval $0\le r\le97/2000$, $v(r)<1$.
The convergent series $1/m=\sum_{k\ge0}v^k$ and all numerators
above have nonnegative coefficients. Thus $C,C',C''\ge0$ there.
For $3\le y\le Y_0$ and $D\ge270$,
\[
 \partial_y^2P_D
 =2C+\frac{4y}{D}C'+\frac{y^2}{D^2}C''
                      +\frac{1000000}{3}\frac{y^3}{D^3}>0,
\]
\[
 \partial_DP_D=-\frac{y^3}{D^2}C'
                    -50000\frac{y^5}{D^4}-(y-2)<0.
\]
All occurrences of $C$ and its derivatives have argument $y/D$.
Direct rational substitution gives
\[
 P_{270}(3)=
 -\frac{343077017913365252920829}{10130931628072339534650}<0,
 \qquad P_{270}(Y_0)<-197.
\]
Convexity bounds $P_{270}$ above by its endpoint chord.
Monotonicity in $D$ then contradicts
\eqref{mf:coarse-bootstrap} if $y\ge3$.

Now $y<3$, hence $q<1/90<7/625$.
The fine quartic row gives
$\Re K'(b)/H^2<(9/4)q^2+13q\mu$.
It follows that
$0<(9/4)y^2+13y-D(y-2)$.
At $D=270$ this convex quadratic has endpoint values
$-1719/64$ and $-843/4$ at $y=9/4$ and $y=3$.
It decreases in $D$ for every $y>2$. Thus $y\ge9/4$ is
impossible. All the rational endpoint inequalities, including
the two companion rows, are checked in
\cert{quartic/verify_d270.py}.
\end{proof}

\subsection{A direct finite-product obstruction}

The following criterion applies to a closed four-dimensional parameter
cell. It uses the original-root factorization in the right half-plane.

\begin{lemma}\label{mf:cubic-moment}
Suppose $\Re w_j\ge0$, $|w_j|\le\mu$, $\sum w_j=\sigma=e^{i\phi}$,
and $\sum|w_j|^2=q$, with $0\le\phi\le\pi/2$.
Put $T=-\sum w_j^2=q(1-\alpha+i\beta)$.
Then
\[
 \alpha^2+\beta^2\le(2\mu\cos\phi/q)^2,\qquad
 (1-\alpha)^2+\beta^2\le1.
\]
If $c=\cos\phi>0$, put $A=q\alpha/2$, $C=-q\beta/2$.
For any upper bound $\mu_h\ge\mu$, the third moment
$P_3=\sum w_j^3$ satisfies
\begin{gather}
 \frac{4A^2}{c}-3\mu_h^2c\le\Re P_3
 \le\min\left(\mu_hA,\mu_h^2c-\frac{4C^2}{c}\right),
 \label{mf:cubic-real}\\
 \left|\Im P_3-\frac{3AC}{c}\right|
 \le\frac32\left(\mu_h^2c-\frac{A^2+C^2}{c}\right)
                         +\mu_h(q-A).
 \label{mf:cubic-imag}
\end{gather}
\end{lemma}
\begin{proof}
Write $w_j=x_j+iy_j$. The identity
$T-q=-2\sum x_jw_j$, weighted Cauchy--Schwarz, and
$\sum x_j=c$ give $|T-q|\le2\mu c$; also $|T|\le q$.
Here $A=\sum x_j^2$, $C=\sum x_jy_j$.
Give the points $(x_j,y_j)$ probability weights $x_j/c$.
Their mean is $(A,C)/c$, and the trace of their covariance
matrix is at most $\mu_h^2-(A^2+C^2)/c^2$.
Now
$\Re P_3=4\sum x_j^3-3\sum x_j|w_j|^2$
gives the lower bound, while
$\Re P_3\le\sum x_j^3\le\mu_h A$ and
$\Re P_3=\sum x_j|w_j|^2-4\sum x_jy_j^2$
give the two upper bounds.
Finally $\Im P_3=3\sum x_j^2y_j-\sum y_j^3$.
The covariance bound controls the first sum about $AC/c$,
and $|\sum y_j^3|\le\mu_h\sum y_j^2=\mu_h(q-A)$
controls the second. Thus points with $x_j=0$ are included.
\end{proof}

\begin{lemma}\label{mf:point-obstruction}
Let $b$ be an actual simple rightmost zero, assume $H=|S|>0$,
and use \eqref{pr:normalized-moments}. Conjugate so that
$\sigma=e^{i\phi}$ with $0\le\phi\le\pi/2$.
Define the actual normalized primitive
\[
 J(w)=\frac{S\,p(b+w/S)}{p'(b)},\quad
 \theta=\frac{-\sum_jw_j^2}{\sigma^2},\quad
 \kappa=\frac{\sum_jw_j^3}{\sigma^3}.
\]
Suppose a closed parameter cell has bounds $q\le q_h$,
$\mu\le\mu_h$, $\phi_l\le\phi\le\phi_h$.
At a nonzero test point $w$, put $W=|w|$ and assume $\mu_hW<1$.
For recorded exact coefficients $\theta_0,\kappa_0$, let
$|\theta-\theta_0|\le h_2$, $|\kappa-\kappa_0|\le h_3$
hold on the whole cell. Put
\[
 B(t)=\sum_{k=4}^8\frac{q_h\mu_h^{k-2}W^k}{k}t^k
       +\frac{q_h\mu_h^7W^9}{9(1-\mu_hW)}t^9,
\]
\[
 J_0=w\int_0^1e^{tw+\theta_0w^2t^2/2+\kappa_0w^3t^3/3}\,dt,
 \qquad g_0=e^{w+\theta_0w^2/2+\kappa_0w^3/3},
\]
\[
 E_J=W\int_0^1 e^{A_0t+B_0t^2+C_0t^3}
     \left(e^{h_2W^2t^2/2+h_3W^3t^3/3+B(t)}-1\right)\,dt,
\]
\[
 E_g=|g_0|\left(e^{h_2W^2/2+h_3W^3/3+B(1)}-1\right),
 \quad
 (A_0,B_0,C_0)=
 \bigl(\Re w,\Re(\theta_0w^2/2),\Re(\kappa_0w^3/3)\bigr).
\]
Write $\phi_m=(\phi_l+\phi_h)/2$,
$\eta_\phi=2\sin((\phi_h-\phi_l)/4)$, $r=|J_0|$, $g=|g_0|$,
and
\[
 P_-=\Re(we^{-i\phi_m})-\eta_\phi W,\quad
 C_+=\Re(e^{i\phi_m}g_0\overline{J_0})+\eta_\phi gr.
\]
If
\begin{equation}\label{mf:point-predicates}
 P_->0,\qquad r>E_J,\qquad
 \frac{C_++gE_J}{r^2-E_J^2}
       +\frac{E_g}{r-E_J}-\frac{P_-}{W^2}<0,
\end{equation}
no actual configuration belongs to that parameter cell.
For a quadratic model, set $\kappa_0=h_3=0$ and begin the
remainder sum at $k=3$; alternatively use
$B(t)=q_h\mu_hW^3t^3/[3(1-\mu_hW)]$.
\end{lemma}

\begin{proof}
For $\Re(w/\sigma)>0$, the original-root factorization
and rightmostness imply
\begin{equation}\label{mf:actual-positivity}
 \Re\{\sigma(J'(w)/J(w)-1/w)\}\ge0.
\end{equation}
Indeed, after multiplication by $H$, the expression is
$p'(z)/p(z)-1/(z-b)=\sum_{u\ne b}(z-u)^{-1}$ at
$z=b+w/S$. Each summand has positive real part.
The critical-point product gives
$J'(w)=\exp(w+\theta w^2/2+\kappa w^3/3+R_4(w))$.
Lemma~\ref{pr:primitive}, with each order through eight kept
separately, gives $|R_4(tw)|\le B(t)$, including every order
at least nine. Integrating the exponential
perturbation proves $|J-J_0|\le E_J$, $|J'-g_0|\le E_g$.

If $|Z-v|\le E<|v|$, inversion maps that disk exactly to
\[
 \left|\frac1Z-\frac{\overline v}{|v|^2-E^2}\right|
 \le\frac{E}{|v|^2-E^2}.
\]
This follows by multiplying the inequality by
$|Z|(|v|^2-E^2)$ and expanding squares.
Apply it to $J$, use the derivative error and
$|e^{i\phi}-e^{i\phi_m}|\le\eta_\phi$, and then subtract
$\Re(\sigma/w)\ge P_-/W^2$.
The resulting upper bound for \eqref{mf:actual-positivity}
is exactly the last expression in \eqref{mf:point-predicates}.
The denominators are positive because $r>E_J$.
\end{proof}

We specify the coefficient enclosures used in this criterion.
For a cell in $(q,\phi,\alpha,\beta)$ take exact midpoints
$q_m,\phi_m,\alpha_m,\beta_m$, set
$C_m=1-\alpha_m+i\beta_m$,
$d_q=(q_h-q_l)/2$, and
$d_C=\sqrt{((\alpha_h-\alpha_l)/2)^2+
                ((\beta_h-\beta_l)/2)^2}$.
Then a valid quadratic radius is
\[
 h_2=d_q(|C_m|+d_C)+q_md_C
       +2q_m|C_m|\sin((\phi_h-\phi_l)/2)
       +|q_mC_me^{-2i\phi_m}-\theta_0|.
\]
Indeed,
\begin{align*}
 &q(1-\alpha+i\beta)e^{-2i\phi}-q_mC_me^{-2i\phi_m}\\
 &\quad=(q-q_m)(1-\alpha+i\beta)e^{-2i\phi}
       +q_m\bigl(\alpha_m-\alpha+i(\beta-\beta_m)\bigr)e^{-2i\phi}\\
 &\qquad\quad+q_mC_m(e^{-2i\phi}-e^{-2i\phi_m}).
\end{align*}
Apply the triangle inequality and
$|e^{-2i\phi}-e^{-2i\phi_m}|
 \le2\sin((\phi_h-\phi_l)/2)$.
For the cubic radius put $c_l=\cos\phi_h$, $c_h=\cos\phi_l$,
$A_l=q_l\alpha_l/2$, $A_h=q_h\alpha_h/2$, and let $C_l,C_h$
be the extrema of $-q\beta/2$ over its four endpoint pairs.
Division is used only when $c_l>0$. Put
$C_*^2=\operatorname{dist}(0,[C_l,C_h])^2$ and
\[
 R_l=4A_l^2/c_h-3\mu_h^2c_h,\qquad
 R_h=\min(\mu_hA_h,\mu_h^2c_h-4C_*^2/c_h).
\]
Let $I_l^0,I_h^0$ be the minimum and maximum of $3AC/c$
over the eight endpoint triples. Set
\[
 V_*=\max(0,\mu_h^2c_h-(A_l^2+C_*^2)/c_h),\qquad
 R_I=\tfrac32V_*+\mu_h(q_h-A_l).
\]
Equations \eqref{mf:cubic-real}--\eqref{mf:cubic-imag} enclose
$P_3$ in $[R_l,R_h]+i[I_l^0-R_I,I_h^0+R_I]$.
Subdivide its entire imaginary interval into $m$ equal closed
strips. For a strip with midpoint $P_m$ and midpoint-to-corner
radius $R_m$, use
\[
 h_3=R_m+2|P_m|\sin(3(\phi_h-\phi_l)/4)
                 +|P_me^{-3i\phi_m}-\kappa_0|.
\]
Every quantity is enclosed in the indicated outward direction.
Every strip is retained, with its endpoints. The values
$m=4,8$ suffice for the two cubic certificates below.

There are three further exclusions for a closed parameter cell.
Let $y_*$ be the threshold $3$ or $9/4$ in the contrary assumption
$q/\mu\ge y_*$, and let $\mu_0$ be the entrance cap, so that
$\mu\le\min(\mu_0,q/y_*)$.
Assume $q_l>0$, $0\le\phi_l\le\phi_h\le\pi/2$ and
$0\le\alpha_l\le\alpha_h<1$, and put
$b_*^2=\operatorname{dist}(0,[\beta_l,\beta_h])^2$.
On the branch $q\le y_*\mu_0$, use
$R_{\mathrm{mom}}=2\cos\phi_l/y_*$ and $\mu_h=q_h/y_*$;
on the branch $q\ge y_*\mu_0$, use
$R_{\mathrm{mom}}=2\mu_0\cos\phi_l/q_l$ and $\mu_h=\mu_0$.
The cell is excluded if any of
\begin{equation}\label{mf:cell-exclusions}
 \alpha_l^2+b_*^2>R_{\mathrm{mom}}^2,\qquad
 (1-\alpha_h)^2+b_*^2>1,\qquad
 q_l>\mathcal E(\mu_h,\phi_l)
\end{equation}
holds, where
\[
 \mathcal E(\mu,\phi)
 =K_0\mu^2+A_0\mu\cos\phi
       +\gamma\cos^2\phi-\tfrac12\sin^2\phi.
\]
Here $K_0,A_0,\gamma$ are the constants in Proposition~\ref{mf:envelope}.
The first two tests use lower bounds for the squared distances in
the two moment disks. For the third, \eqref{mf:exact-envelope} gives
$q\le\mathcal E(\mu,\phi)$, while
\[
 \partial_\mu\mathcal E=2K_0\mu+A_0\cos\phi\ge0,\qquad
 \partial_\phi\mathcal E
 =-\sin\phi\{A_0\mu+(2\gamma+1)\cos\phi\}\le0.
\]
All three tests therefore exclude the whole cell; their strict
inequalities are checked with outward enclosures.

\begin{proposition}\label{mf:D30}
If $D=Hd\ge30$, then $Qd/H<3$.
\end{proposition}
\begin{proof}
Proposition~\ref{mf:D270} treats $D\ge270$. Suppose
$30\le D\le270$ and $q/\mu\ge3$.
Proposition~\ref{mf:affine}, in its consequence for $D\ge7/2$,
and Lemma~\ref{mf:cubic-moment} put the contrary domain in
\[
 q\in[1/90,403/900],\quad
 \mu\le\min(1/30,q/3),\quad
 0\le\phi\le\pi/2,\quad
 0\le\alpha\le2/3,\quad|\beta|\le2/3,
\]
with the two moment-disk inequalities and the exact matrix
envelope. Divide the $q$ interval at $1/10$.
The fixed dyadic trees in
\cert{fields/tang_cubic_family_D30_y3.json}
have $9613$ and $16784$ closed leaves. A leaf is excluded by
\eqref{mf:cell-exclusions} or Lemma~\ref{mf:point-obstruction}.
The certificate has $3237,239,390$ leaves of the first three
types, $3423$ quadratic-model leaves, $8$ single cubic-model
leaves, and $19100$ complete cubic-strip families.
The latter have $79392$ members. Thus there are $82823$
point obstructions in all.

The verifier reconstructs each leaf from the exact dyadic
path; both halves of every subdivision are closed. Bottom-up
sibling coalescence reduces each tree to its root,
rejecting duplicate leaves, missing siblings, ancestor conflicts,
and inconsistent splitting coordinates. All inequalities and
all the point predicates are checked on the entire leaves by
\cert{fields/tang_verify_cubic_family_cover_D30_frozen.py}.
The source file records the exact test points and coefficient
centers. Stored finite decimals defining those centers are
interpreted as their exact binary rational values; their
mismatch from the mathematical centers is included in the
radii above. The integrands are entire, with integration remainders
bounded as in Section~\ref{sec:certificates}.
The certificate excludes the stated closed domain.
Together with $D\ge270$ this proves the proposition.
\end{proof}

\begin{proposition}\label{mf:D25}
If $D=Hd\ge25$, then $Qd/H<3$.
\end{proposition}
\begin{proof}
Only $25\le D\le30$ remains after Proposition~\ref{mf:D30}.
Under the contrary ratio $q/\mu\ge3$, the affine estimate gives
\[
 1/10\le q\le59/125,\quad \mu\le\min(1/25,q/3),\quad
 0\le\alpha\le2/3,\quad|\beta|\le2/3.
\]
In this range $\phi<\pi/4$. Indeed if $\phi\ge\pi/4$,
\eqref{mf:coarse-envelope}, $\cos\phi<71/100$, and
$q\ge3\mu$ give
\[
 0<\frac{21}{5}\mu^2+\frac{181}{200}\mu-\frac18
 \le-\frac{513}{6250}<0.
\]
Use the closed $\phi$ interval $[0,\pi/4]$ and divide
the $q$ interval at $3/25$.
The two dyadic trees in
\cert{fields/tang_D25_cubic_family_cover.json}
have $15303$ and $35311$ leaves. The exclusions consist of
$4752$ moment-disk, $393$ quadratic-disk, and $219$ matrix
leaves, using \eqref{mf:cell-exclusions}.
There are $2246$ quadratic-model leaves, $15$ single
cubic-model leaves, and $42989$ complete strip families;
the families contain $184216$ members.
Every one of the $186477$ point obstructions satisfies
\eqref{mf:point-predicates}.
The checker \cert{fields/tang_verify_D25_cubic_family_frozen.py}
uses the preceding finite-product and covariance formulas and
reconstructs the two closed trees.
Thus it excludes this domain. Its common boundary
with $D\ge30$ is included.
\end{proof}

\begin{proposition}\label{mf:D100}
If $D=Hd\ge100$, then $Qd/H<9/4$.
\end{proposition}
\begin{proof}
By Proposition~\ref{mf:D270}, it remains to consider
$100\le D\le270$ and $q/\mu\ge9/4$.
The contrary domain is contained in
\[
 q\in[1/120,1313/10000],\quad
 \mu\le\min(1/100,4q/9),\quad
 0\le\phi\le\pi/2,\quad
 0\le\alpha\le8/9,\quad|\beta|\le8/9,
\]
with the same two moment disks and matrix envelope.
Split the $q$ interval at $9/400$. The two trees in
\cert{fields/point_D100_y225.json} have $5200$ and $4545$ leaves.
There are $1507$ moment-disk exclusions, $57$ quadratic-disk
exclusions, and $21$ matrix exclusions from
\eqref{mf:cell-exclusions}. On each of the
remaining $8160$ leaves the quadratic version of
Lemma~\ref{mf:point-obstruction} applies, with
$B(t)=q_h\mu_hW^3t^3/[3(1-\mu_hW)]$.
The coefficient radius is the full-cell $h_2$ displayed above.
The checker \cert{fields/tang_independent_D100.py} reconstructs
both trees, all coefficient radii, and every strict
point obstruction. It evaluates the equivalent entire integral
$J_0/g_0=w\int_0^1e^{(t-1)w+(t^2-1)\theta_0w^2/2}\,dt$;
this follows by dividing the defining integral by $g_0$.
Thus every point of the contrary domain is excluded, which
proves the proposition.
\end{proof}
